\documentclass[11pt,a4paper]{article}
\usepackage[margin=3cm]{geometry}
\usepackage{amsmath,amssymb,amsthm,mathrsfs}
\usepackage[T1]{fontenc}
\usepackage{microtype}
\usepackage{booktabs}
\usepackage[colorlinks=true,linkcolor=blue,citecolor=blue]{hyperref}

\theoremstyle{plain}
\newtheorem{theorem}{Theorem}[section]
\newtheorem{proposition}[theorem]{Proposition}
\newtheorem{lemma}[theorem]{Lemma}
\newtheorem{corollary}[theorem]{Corollary}
\theoremstyle{definition}
\newtheorem{remark}[theorem]{Remark}
\newtheorem{hypothesis}[theorem]{Hypothesis}

\newcommand{\A}{\mathfrak{A}}
\newcommand{\R}{\mathbb{R}}
\newcommand{\Om}{\Omega}
\newcommand{\ip}[2]{\langle #1,#2\rangle}
\DeclareMathOperator{\dist}{dist}
\DeclareMathOperator{\supp}{supp}
\DeclareMathOperator{\spec}{spec}

\title{On (H6): the local second moment of the interaction density\\
\large what the gap cannot do, and how much less than (H6) suffices}
\author{}
\date{}

\begin{document}
\maketitle

\begin{abstract}
(H6) asks for $\sup_{g}\|V_B\Om_g\|<\infty$, $V_B$ a local Wick polynomial. \textbf{I do not have
a proof.} This note contributes three things.

\smallskip
\noindent\textbf{(1) A negative result: the gap is the wrong tool, for an ultraviolet reason.}
The two natural gap-based routes are \emph{vacuous}. Both pass through quantities that are
UV divergent while $\|V_B\Om_g\|$ itself is finite. The $f$-sum rule gives
$\ip{V_B\Om_g}{H(g)V_B\Om_g}=\tfrac12\,\omega_g\bigl([V_B,[H_0,V_B]]\bigr)$ --- note the double
commutator is \emph{$g$-independent}, because $[V(g),V_B]=0$ --- and equals
$\tfrac12\omega_g\bigl(\int h^2(:\!P'(\varphi)\!:)^2\bigr)$, a Wick square at coincident points.
In the free case it diverges logarithmically; verified numerically (\S\ref{sec:neg}).

\smallskip
\noindent\textbf{(2) A positive result: (H6) can be weakened to (H6$'$), uniform
$L^{1+\delta}$.} Truncating $V_B$ at height $M$ and using that the \emph{exact light cone} makes
the truncation tail \emph{commute} with $\alpha_t(A)$ for $|t|<R$, one needs only a first moment
of the tail. This costs a factor $\delta/(1+\delta)$ in the clustering rate
(Theorem~\ref{thm:trunc}). A first moment alone --- which \emph{is} available, uniformly, from
Guerra--Rosen--Simon --- is not enough; uniform integrability is exactly what is missing.

\smallskip
\noindent\textbf{(3) A sharp identification of what remains.} Via H\"older and Nelson
hypercontractivity, (H6$'$) follows from a uniform local $L^q$ bound on the ground-state density
for \emph{some} $q>1$, i.e.\ from Problem~1 of \cite{GRS75} verbatim; and via localised $N_\tau$
estimates, (H6) follows from the $\deg P$-th moment of the local number operator, strictly
stronger than (H2)-full's first moment (\S\ref{sec:red}).
\end{abstract}

\section{Setting; (H6) in $Q$-space}

Notation as in \cite{M0,M1,M3}. Let $h\in C_c^\infty(B)$, $|h|\le1$, and
\[
  V_B:=\int_B h(x)\,{:}P(\varphi_0(x)){:}\,dx ,
\]
a symmetric operator affiliated to $\A(B)$. In the $Q$-space representation $V_B$ is a
\emph{multiplication} operator and $d\mu_g:=\Om_g^2\,d\varphi_0$ is a probability measure, so
\begin{equation}\label{eq:qspace}
  \|V_B\Om_g\|^2=\mathbb E_{\mu_g}\bigl[V_B^2\bigr],
\end{equation}
and (H6) reads: \emph{the second moment of a fixed local random variable is bounded uniformly in
the cutoff}. For each fixed $g$ it is finite, since $\Om_g\in L^p(Q,d\varphi_0)$ for all
$p<\infty$ and $V_B\in L^p$ for all $p<\infty$ by Nelson's hypercontractive bound; only the
uniformity in $g$ is at stake.

Two facts are available \emph{uniformly} in $g$ and will be used repeatedly:
\begin{enumerate}
\item[\textbf{(F$_1$)}] \emph{Uniform first moment.} By \cite[Thm.~2]{GRS72}, $\pm V_B\le\varepsilon
      H(g)+c(B)$ with $c$ independent of $g$; taking $\omega_g$ and using $H(g)\Om_g=0$,
      $|\omega_g(V_B)|\le c(B)$, and moreover
      $\|(H(g)+c)^{-1/2}V_B(H(g)+c)^{-1/2}\|\le1+\varepsilon$.
\item[\textbf{(F$_2$)}] \emph{Uniform local excitation energy.} $\|H(g)^{1/2}A\Om_g\|^2
      =-i\,\omega_g(A^*\delta(A))$ by \cite[Lem.~2.1]{M0}, bounded uniformly in $g$ for $A$ in the
      class covered by the $\varphi$-bound of \cite{GJIV}.
\end{enumerate}

\section{The gap cannot prove (H6)}\label{sec:neg}

\begin{proposition}[$f$-sum rule; the double commutator is $g$-independent]\label{prop:fsum}
For symmetric $X$ with $X\Om_g\in D(H(g)^{1/2})$,
\[
  \ip{X\Om_g}{H(g)X\Om_g}=\tfrac12\,\omega_g\bigl([X,[H(g),X]]\bigr).
\]
If $X=V_B$ then, because $V_B$ and the interaction $V(g)$ are both functions of $\varphi_0$ alone
and hence commute,
\[
  [V_B,[H(g),V_B]]=[V_B,[H_0,V_B]]=\int h(y)^2\bigl({:}P'(\varphi_0(y)){:}\bigr)^2dy .
\]
\end{proposition}

\begin{proof}
$H(g)\Om_g=0$ gives $\omega_g([X,[H,X]])=2\ip{X\Om_g}{HX\Om_g}$ directly. For the second claim,
$[V(g),V_B]=0$ and $[E(g),V_B]=0$, so only $H_0$ contributes; only the $\tfrac12\int{:}\pi^2{:}$
term of $H_0$ fails to commute with a function of $\varphi_0$. With
$[\,V_B,\pi(y)]=i\,h(y){:}P'(\varphi_0(y)){:}$ one computes
$[H_0,V_B]=-\tfrac i2\int h(y)\bigl(\pi(y){:}P'{:}+{:}P'{:}\pi(y)\bigr)dy$ and then the stated
identity.
\end{proof}

Combining Proposition~\ref{prop:fsum} with the gap $H(g)\ge\gamma(1-P_g)$ gives the tempting
\begin{equation}\label{eq:tempting}
  \|V_B\Om_g\|^2\ \le\ \omega_g(V_B)^2+\frac{1}{2\gamma}\,
  \omega_g\Bigl(\int h^2\bigl({:}P'(\varphi_0){:}\bigr)^2\Bigr),
\end{equation}
whose first term is uniformly bounded by (F$_1$) and whose second term is $g$-independent. This
looks like a proof of (H6). It is not.

\begin{proposition}[the bound is vacuous]\label{prop:vacuous}
The right-hand side of \eqref{eq:tempting} is $+\infty$. Already for $P(\xi)=\xi^2$ and the Fock
vacuum,
\[
  \ip{V_B\Om_0}{H_0V_B\Om_0}\;=\;c\!\int\!\!\int
  |\hat h(k_1+k_2)|^2\,\frac{\mu(k_1)+\mu(k_2)}{\mu(k_1)\mu(k_2)}\,dk_1dk_2\;=\;+\infty ,
\]
logarithmically divergent, while $\|V_B\Om_0\|^2=c\iint|\hat h(k_1+k_2)|^2(\mu_1\mu_2)^{-1}<\infty$.
\end{proposition}

\begin{proof}
$V_B\Om_0$ is the two-particle vector with kernel $\propto\hat h(k_1+k_2)(\mu_1\mu_2)^{-1/2}$.
Writing $K=k_1+k_2$, the norm integral is $\int|\hat h(K)|^2\!\int dk\,(\mu(k)\mu(K-k))^{-1}$,
convergent since the inner integrand is $O(k^{-2})$; the energy integral is
$\int|\hat h(K)|^2\!\int dk\,(\mu(k)^{-1}+\mu(K-k)^{-1})$, whose inner integrand is $O(|k|^{-1})$.
\end{proof}

\emph{Numerical confirmation} ($m=1$, $|\hat h(K)|^2=e^{-K^2/2}$, symmetric cutoff $|k_i|\le\Lambda$):

\begin{center}
\begin{tabular}{rrr}
\toprule
$\Lambda$ & $\|V_B\Om_0\|^2$ & $\ip{V_B\Om_0}{H_0V_B\Om_0}$\\
\midrule
$20$  & $7.2493$ & $36.79$\\
$80$  & $7.4419$ & $50.84$\\
$320$ & $7.4892$ & $64.78$\\
$640$ & $7.4970$ & $71.73$\\
\bottomrule
\end{tabular}
\end{center}
The norm saturates; the energy increases by $\approx6.96$ per doubling of $\Lambda$ --- exactly
logarithmic.

\begin{remark}[the resolvent route fails identically]
From $H(g)V_B\Om_g=[H_0,V_B]\Om_g$ one gets
$\|(1-P_g)V_B\Om_g\|\le\gamma^{-1}\|[H_0,V_B]\Om_g\|$; but $\|[H_0,V_B]\Om_0\|=\infty$ by the same
computation. Vacuous again.
\end{remark}

\begin{remark}[diagnosis]\label{rem:diag}
$\|V_B\Om_g\|^2$ is \emph{ultraviolet}-sensitive: in $Q$-space it is
$\mathbb E_{\mu_g}[V_B^2]$, and the Wick reordering
${:}P(\varphi(x)){:}{:}P(\varphi(y)){:}=\sum_k\frac{C(x-y)^k}{k!}
{:}P^{(k)}(\varphi(x))P^{(k)}(\varphi(y)){:}$ makes it depend on the short-distance behaviour
$C(x-y)\sim-\tfrac1{2\pi}\log|x-y|$. The spectral gap is an \emph{infrared} object. Any argument
routed through $H(g)^{-1}(1-P_g)$ or through a double commutator discards the Wick-ordering
structure and produces a divergent majorant. \textbf{The tools that can work are
hypercontractivity, $N_\tau$ estimates and local $L^p$ bounds --- not the gap.}
\end{remark}

\section{(H6) can be weakened to (H6$'$)}\label{sec:trunc}

\begin{hypothesis}[(H6$'$)]\label{hyp:H6p}
There are $\delta>0$ and $c_\delta(B)<\infty$ with
$\displaystyle\sup_{g\in\mathcal G}\ \mathbb E_{\mu_g}\bigl[|V_B|^{1+\delta}\bigr]\le c_\delta(B)$.
\end{hypothesis}

The point of the next theorem is that the \emph{exact light cone} makes the truncation tail
commute with the evolved observable, so that only a \emph{first} moment of the tail is needed
where a Cauchy--Schwarz step would have demanded a second.

\begin{theorem}[truncation]\label{thm:trunc}
Let $A\in\A(O)$, $R:=\dist(O,B)>0$, and let $g$ satisfy $d(g,O)>R$ and
$\spec H(g)\subseteq\{0\}\cup[\gamma,\infty)$. Put
$\Psi(t):=\omega_g\bigl(V_B\alpha_t(A)\bigr)-\omega_g(V_B)\omega_g(A)$. Assume (H6$'$). Then for
all $|t|<R$,
\[
  |\Psi(t)|\ \le\ 2\bigl(1+c_\delta(B)\bigr)\|A\|\;e^{-\gamma'\,(R-|t|)},
  \qquad \gamma':=\frac{\gamma\,\delta}{2(1+\delta)} .
\]
\end{theorem}

\begin{proof}
Fix $M>0$ and set $V^{(M)}:=V_B\mathbf 1_{\{|V_B|\le M\}}\in\A(B)$, a bounded multiplication
operator with $\|V^{(M)}\|\le M$, and $Y:=V_B-V^{(M)}=V_B\mathbf 1_{\{|V_B|>M\}}$. Write
$Y=Y_+-Y_-$ with $Y_\pm\ge0$, both affiliated to $\A(B)$.

\emph{Step 1 (the tail costs only a first moment).} Let $|t|<R$. By the exact light cone
\cite[Fact~1.3(i)]{M0}, $\alpha_t(A)\in\A(O_{|t|})$ and $O_{|t|}\cap B=\emptyset$, so
$\alpha_t(A)$ commutes with $Y_\pm$ and hence with $Y_\pm^{1/2}$. Therefore
\[
  \bigl|\ip{Y_\pm\Om_g}{\alpha_t(A)\Om_g}\bigr|
  =\bigl|\ip{Y_\pm^{1/2}\Om_g}{\alpha_t(A)\,Y_\pm^{1/2}\Om_g}\bigr|
  \le\|A\|\,\bigl\|Y_\pm^{1/2}\Om_g\bigr\|^2=\|A\|\,\mathbb E_{\mu_g}[Y_\pm].
\]
Adding the two signs and the disconnected term,
\[
  \bigl|\Psi(t)-\Psi^{(M)}(t)\bigr|\ \le\ 2\|A\|\;\mathbb E_{\mu_g}\bigl[|V_B|\mathbf 1_{\{|V_B|>M\}}\bigr]
  \ \le\ 2\|A\|\,c_\delta(B)\,M^{-\delta},
\]
the last step by $|V_B|\mathbf 1_{\{|V_B|>M\}}\le|V_B|^{1+\delta}M^{-\delta}$ and (H6$'$). Here
$\Psi^{(M)}$ is $\Psi$ with $V_B$ replaced by $V^{(M)}$.

\emph{Step 2 (clustering for the bounded part).} $V^{(M)}\in\A(B)$ and $A\in\A(O)$ with
$\dist(O,B)=R$; the argument of \cite[Thm.~2.3]{M3} applied at separation $R-|t|$ gives
$|\Psi^{(M)}(t)|\le2M\|A\|e^{-\gamma(R-|t|)/2}$.

\emph{Step 3 (optimise).} Choose $M:=\exp\bigl(\tfrac{\gamma(R-|t|)}{2(1+\delta)}\bigr)$. Then
$Me^{-\gamma(R-|t|)/2}=M^{-\delta}=e^{-\gamma'(R-|t|)}$, and adding the two bounds gives the claim.
\end{proof}

\begin{remark}[a first moment is \emph{not} enough]
(F$_1$) supplies $\sup_g\mathbb E_{\mu_g}[|V_B|]<\infty$ uniformly. That alone does not bound
$\mathbb E_{\mu_g}[|V_B|\mathbf 1_{\{|V_B|>M\}}]$ uniformly in $g$: what is required is
\emph{uniform integrability} of the family $\{V_B\}$ under $\{\mu_g\}$, and (H6$'$) is precisely
the standard sufficient condition for it. So the gap between what is known and what is needed is
exactly ``$L^1$ uniformly'' versus ``$L^{1+\delta}$ uniformly''.
\end{remark}

\begin{lemma}[the far-time regime needs nothing new]\label{lem:fartime}
For all $t\in\R$, $|\Psi(t)|\le K(A):=\sqrt c\,(1+\varepsilon)\bigl(\|(H(g)+c)^{1/2}A\Om_g\|
+\sqrt c\,\|A\|\bigr)$, uniformly in $g$ and $t$.
\end{lemma}

\begin{proof}
$\Psi(t)=\ip{\Om_g}{V_B\bigl(\alpha_t(A)-\omega_g(A)\bigr)\Om_g}$; insert
$(H+c)^{\pm1/2}$ on both sides of $V_B$ and use (F$_1$), $\|(H+c)^{1/2}\Om_g\|=\sqrt c$, and
$\|(H+c)^{1/2}\alpha_t(A)\Om_g\|=\|(H+c)^{1/2}e^{itH}A\Om_g\|=\|(H+c)^{1/2}A\Om_g\|$, which is
$t$-independent and uniformly bounded by (F$_2$).
\end{proof}

Theorem~\ref{thm:trunc} and Lemma~\ref{lem:fartime} together supply exactly the two regimes the
M4 estimate needs, with (H6) replaced by (H6$'$) and no other new input.

\section{What (H6$'$) still requires}\label{sec:red}

\begin{proposition}\label{prop:red}
Let $\mu_g^B,\varphi_0^B$ denote the restrictions to the $\sigma$-algebra generated by the field in
$B$. If for some $q>1$
\begin{equation}\label{eq:P1}
  \sup_{g\in\mathcal G}\ \Bigl\|\frac{d\mu_g^B}{d\varphi_0^B}\Bigr\|_{L^q(\varphi_0^B)}<\infty ,
\end{equation}
then \textup{(H6$'$)} holds for every $\delta>0$; and if \eqref{eq:P1} holds for some $q>2$ then
\textup{(H6)} holds.
\end{proposition}

\begin{proof}
$V_B$ is $\A(B)$-measurable, so $\mathbb E_{\mu_g}[|V_B|^{r}]=\mathbb E_{\mu_g^B}[|V_B|^{r}]$.
By H\"older with $1/p+1/q=1$,
$\mathbb E_{\mu_g^B}[|V_B|^{r}]\le\||V_B|^{r}\|_{L^p(\varphi_0^B)}\,\|d\mu_g^B/d\varphi_0^B\|_{L^q}$,
and $\||V_B|^r\|_{L^p(\varphi_0)}=\|V_B\|_{L^{rp}(\varphi_0)}^r<\infty$ for every $rp<\infty$ by
Nelson's hypercontractive bound, with a constant depending only on $B,P$ and not on $g$.
\end{proof}

\begin{remark}[this is GRS Problem 1, verbatim]
\eqref{eq:P1} is Problem~1 of \cite[p.~125]{GRS75}: ``prove that for some fixed $p>1$,
$\|\nu_{\Lambda\Lambda'}\|_p$ is bounded independently of $\Lambda'$'', with the conjecture
$\log\|\nu_{\Lambda\Lambda'}\|_p=O(|\Lambda|)$. GRS observe that such a bound would re-prove the
locally Fock property and give the infinite-volume equal-time vacuum expectation values. So
(H6$'$) is not a new problem; it is a fifty-year-old one, and Theorem~\ref{thm:trunc} shows the
weakest form of it ($q>1$, arbitrarily close to $1$) already suffices for M4.
\end{remark}

\begin{remark}[the $N_\tau$ route, for comparison]
The localised $N_\tau$ estimates of \cite[(3.2.9), (4.19)]{GJIII} give
$\|V_B\Om_g\|\le\text{const}\cdot\|w_B\|\cdot\|(N_{B'}+I)^{n/2}\Om_g\|$ for a degree-$n$ Wick
monomial, so (H6) follows from $\sup_g\omega_g\bigl(N_{B'}^{\,n}\bigr)<\infty$ with
$n=\deg P$ --- the $\deg P$-th moment of the local number operator, strictly stronger than
(H2)-full's first moment. This route is therefore worse than Proposition~\ref{prop:red}.
\end{remark}

\section{Status and what I would try next}

\begin{itemize}
\item \textbf{Proved:} the gap-based routes are vacuous, for an ultraviolet reason
      (Props.~\ref{prop:fsum}--\ref{prop:vacuous}, Rem.~\ref{rem:diag}); (H6) may be replaced by
      the strictly weaker (H6$'$) at a rate cost $\delta/(1+\delta)$ (Thm.~\ref{thm:trunc}); the
      far-time regime needs no new input (Lem.~\ref{lem:fartime}).
\item \textbf{Not proved:} (H6$'$). It is implied by, and appears to need, a uniform local $L^q$
      bound on the ground-state density with $q>1$ --- GRS Problem~1.
\item \textbf{Most promising input.} Bauerschmidt--Dagallier's log-Sobolev inequality for the
      $\varphi^4_2$ measure is uniform in the volume \emph{and} in the lattice regularisation.
      LSI is stable under marginalisation, so a uniform LSI for the space-cutoff Euclidean measure
      $\nu_g$ descends to its time-zero marginal $\mu_g$; and an LSI yields Gaussian concentration
      for suitable functionals, which is the natural source of uniform $L^{1+\delta}$ bounds on a
      fixed local $V_B$. Whether this can be pushed to \eqref{eq:P1}, or directly to (H6$'$) for
      Wick polynomials rather than Lipschitz functions, is the concrete question. It is already on
      the M5 critical path, so the cost of checking is shared.
\end{itemize}

\section{Two candidate routes to (H6$'$), and a pattern}\label{sec:routes}

\subsection*{Route~1: Gaussian domination of moments --- blocked exactly at Wick powers}

For even $P$ in the Griffiths--Simon class the lattice approximation is a genuine ferromagnet
\cite[\S IV.4]{GRS75}, and there the machinery exists: Sylvester \cite{Syl75} and Shlosman
\cite{Shl86} give the signs of the Ursell functions ($u_{2n}\le0$ for $n\ge2$), and Newman
\cite{New75}, Ellis--Monroe--Newman \cite{EMN76} convert these into \emph{Gaussian upper bounds}
\[
  \langle\varphi_{x_1}\cdots\varphi_{x_{2n}}\rangle\ \le\
  \sum_{\text{pairings}}\ \prod\langle\varphi\varphi\rangle .
\]
If such a bound survived to the continuum, (H6) would follow \emph{outright}, and with the
\emph{free-field} value as the uniform constant:
$\mathbb E_{\mu_g}[V_B^2]\le\mathbb E_{\varphi_0}[V_B^2]=(\deg P)!\iint h\,h\,C^{\deg P}<\infty$,
finite because $C(x-y)\sim-\tfrac1{2\pi}\log|x-y|$ is locally $L^p$ for every $p$ in one dimension.

\textbf{The obstruction is the Wick subtraction, and it is documented.} On the lattice there is no
Wick ordering; in the continuum limit the counterterms subtract $C_\varepsilon(0)\sim\log(1/\varepsilon)$,
and an \emph{inequality} does not survive subtracting a divergent quantity from both sides. This is
exactly the complaint of \cite[\S V.3]{GRS75}: ``the fact that the interaction involves Wick powers
limits the applicability of the correlation inequalities \dots we are able to handle only a single
quadratic power ${:}\varphi^2{:}$''; and the situation has not materially changed.

\emph{The concrete sub-problem}, therefore: does the Gaussian bound survive Wick subtraction, given
$\langle\varphi(x)^2\rangle_g\le C(0)$ (expected for repulsive $P$)? This is a finite combinatorial
question --- matching counterterms on the two sides of the inequality --- not an analytic one, and
is worth an afternoon.

\subsection*{Route~2: chessboard estimates --- and a synergy with M1}

Multiple-reflection (chessboard) estimates in the sense of Fr\"ohlich--Israel--Lieb--Simon convert a
\emph{local} expectation into a universal bound controlled by the \emph{pressure}, and the pressure
is uniformly bounded here: $-MV\le E_g\le0$ \cite[Thm.~5.1]{GJIII} and $\alpha(l)\uparrow\alpha_\infty<\infty$
\cite[Thm.~1]{GRS72}. This route does not care about Wick powers in the way Route~1 does.

\emph{Tension:} reflection positivity requires the cutoff to respect the reflection hyperplanes,
which a generic $g$ does not. \emph{Resolution:} by \cite[Thm.~3.3]{M1} we are free to choose the
cofinal family --- the Chebyshev selection already exploits that freedom --- so one may take
reflection-symmetric or periodic cutoffs. This is a genuine opening and costs nothing structurally.

\subsection*{Dead ends, recorded so they are not repeated}

\begin{enumerate}
\item \emph{Interpolating from the first moment.} $\mathbb E[|V|^{1+\delta}]\le
      \mathbb E[|V|]^{1-\theta}\mathbb E[|V|^{r}]^{\theta}$ needs a higher moment: circular.
\item \emph{\cite[Thm.~2]{GRS72} with $Q=\pm\varepsilon P^2$.} $P+\varepsilon P^2$ is semibounded, so one
      obtains only $\omega_g(\int h{:}P^2(\varphi){:})\ge-c/\varepsilon$ --- a \emph{lower} bound.
      $P-\varepsilon P^2$ is not semibounded, so the theorem gives no upper bound. (And
      $\int h{:}P^2{:}\ne V_B^2$ in any case.)
\item \emph{Euclidean conditioning.} Writing the local density via the Markov property produces an
      uncontrolled boundary functional --- that \emph{is} Problem~1.
\item \emph{The spectral gap.} Dead, for the ultraviolet reason of \S\ref{sec:neg}.
\end{enumerate}

\subsection*{A pattern worth noticing}

Correlation inequalities are available for ${:}\varphi^2{:}$ and no further. That is the same
boundary as M2's solvable model (quadratic $P$), and the same boundary as the FKG/Daubechies
dichotomy of the wavelet programme (ferromagnetic structure at $K=1$, lost for $K\ge2$). Three
independent places where the difficulty is exactly \emph{``positivity is available at the bottom of
the hierarchy and lost above''}. This is evidence that the log-Sobolev / Polchinski route --- which
does not require positivity --- is the structurally correct bet, and it is already the M5 plan.

\begin{thebibliography}{9}
\bibitem{M0} Milestone M0, \texttt{../m0/m0.tex}.
\bibitem{M1} Milestone M1, \texttt{../m1/m1.tex}.
\bibitem{M3} Milestone M3, \texttt{../m3/m3.tex}.
\bibitem{GJIII} J.~Glimm and A.~Jaffe, Acta Math.\ \textbf{125} (1970) 203--261.
\bibitem{GJIV} J.~Glimm and A.~Jaffe, J.\ Math.\ Phys.\ \textbf{13} (1972) 1568--1584.
\bibitem{GRS72} F.~Guerra, L.~Rosen and B.~Simon, Comm.\ Math.\ Phys.\ \textbf{27} (1972) 10--22.
\bibitem{GRS75} F.~Guerra, L.~Rosen and B.~Simon, Ann.\ of Math.\ \textbf{101} (1975) 111--189,
191--259.
\bibitem{Syl75} G.~S.~Sylvester, \emph{Representations and inequalities for Ising model Ursell
functions}, Comm.\ Math.\ Phys.\ \textbf{42} (1975) 209--220.
\bibitem{Shl86} S.~B.~Shlosman, \emph{Signs of the Ising model Ursell functions}, Comm.\ Math.\
Phys.\ \textbf{102} (1986) 679--686.
\bibitem{New75} C.~M.~Newman, \emph{Gaussian correlation inequalities for ferromagnets},
Z.\ Wahrsch.\ Verw.\ Gebiete \textbf{33} (1975) 75--93.
\bibitem{EMN76} R.~S.~Ellis, J.~L.~Monroe and C.~M.~Newman, \emph{The GHS and other correlation
inequalities for a class of even ferromagnets}, Comm.\ Math.\ Phys.\ \textbf{46} (1976) 167--182.
\end{thebibliography}

\end{document}
